\begin{table}[htbp]
\centering
\begin{threeparttable}
\caption{Within-auction share under positive cost affiliation (Gaussian copula)}
\label{tab:affiliation_sensitivity}
\footnotesize
\setlength{\tabcolsep}{4pt}
\begin{tabular}{lrrrr}
\toprule
Within-auction cost correlation $\rho_c$ & 0.0 (IPV) & 0.1 & 0.2 & 0.3 \\
\midrule
Non-pharma within-auction share & 74.5\% & 73.8\% & 72.6\% & 70.9\% \\
Pharma within-auction share & 73.3\% & 72.5\% & 71.0\% & 68.8\% \\
\midrule
$p_{S_3} - p_{S_1}$ non-pharma & $+0.259$ & $+0.255$ & $+0.249$ & $+0.241$ \\
$p_{S_3} - p_{S_1}$ pharma & $+0.308$ & $+0.302$ & $+0.293$ & $+0.281$ \\
\bottomrule
\end{tabular}
\begin{tablenotes}\footnotesize
\item Sensitivity of the within-auction share and total simulated price effect
under positive cost affiliation. Cost draws are sampled from the type-specific
$F_c^k$ marginals coupled by a Gaussian copula with within-auction correlation
$\rho_c \in \{0, 0.1, 0.2, 0.3\}$. $\rho_c = 0$ recovers the IPV baseline.
Affiliation up to $\rho_c = 0.3$ moves the within-auction share by less than
5 percentage points and reduces the total simulated effect by less than 10 percent.
The 73--85 percent range claimed across the cost-distribution-estimator robustness
remains valid under affiliation up to $\rho_c = 0.3$. Higher affiliation levels
($\rho_c \ge 0.5$) are inconsistent with the cross-modality consistency check
(Section~\ref{sec:identification}, Table~\ref{tab:v3_cross_modality}), which
rejects the alternative-format GPV-vs-drop-out divergence implied by strong
common-values contamination.
\end{tablenotes}
\end{threeparttable}
\end{table}
